\section{Decision Intelligence}
\label{sec:intelligence}

The central contribution of the second project phase is a deterministic decision pipeline that operates on canonical billing and allocation outputs rather than on dashboard-local heuristics. Starting from the unified billing mart, the system materializes signals, recommendations, forecast outputs, and lifecycle summaries as explicit artifacts. This makes the intelligence layer auditable: every recommendation can be traced back to a normalized billing condition, a fixed mapping rule, and a persisted score decomposition.

\subsection{Waste Detection}
The waste detector classifies inefficiency into a finite taxonomy rather than relying on a generic overspend flag. The implemented categories are unused commitment, underutilized commitment, idle compute proxy, and zombie resource. Thresholds are loaded from configuration, allowing sensitivity to be tuned without changing code.

\begin{table}[t]
\caption{Implemented waste taxonomy.}
\label{tab:waste-taxonomy}
\centering
\footnotesize
\begin{tabularx}{\columnwidth}{p{0.18\columnwidth} p{0.24\columnwidth} p{0.12\columnwidth} Y}
\toprule
\textbf{Waste type} & \textbf{Signal} & \textbf{Conf.} & \textbf{Main limitation} \\
\midrule
Unused commitment & Explicit commitment waste rows & High & Depends on provider waste rows being exposed correctly \\
Underutilized commitment & Commitment utilization below configured threshold & Medium & Sensitive to threshold choice \\
Idle compute proxy & NEC-per-vCPU proxy below configured floor & Medium & Billing-derived proxy, not runtime telemetry \\
Zombie resource & Near-zero NEC over the observation window & Medium & Low-cost resources may still have latent purpose \\
\bottomrule
\end{tabularx}
\end{table}

The distinction between deterministic and proxy signals is important. Unused commitment is close to deterministic because it is derived from explicit waste fields. By contrast, idle compute is a billing-side approximation of low utilization; the system intentionally does not claim runtime observability it does not have.

\subsection{Structured Reasoning and Recommendation Scoring}
The next layer performs structured reasoning over billing-derived facts. Methodologically, it is better understood as a ranked explanation layer over cost signals, trends, and anomalies than as intervention-based causal inference. Fact types such as commitment waste, idle resources, unattributed cost gaps, and service-category dominance are combined into persisted explanatory narratives.

Recommendations are generated from canonical signals using an explicit signal-to-action map. Idle compute proxy findings map to \texttt{resize\_down}, zombie resource findings map to \texttt{remove\_resource}, governance gaps map to \texttt{enforce\_tags}, and forecast risk signals map to \texttt{review\_forecast\_risk}. Recommendation ranking is deterministic:
\begin{equation}
p = 0.35s + 0.25c + 0.20g + 0.10u + 0.10l
\end{equation}
where $p$ is priority score, $s$ is normalized savings, $c$ is confidence, $g$ is governance severity, $u$ is urgency, and $l$ is a low-risk bonus.

\begin{table}[t]
\caption{Representative signal-to-action mappings.}
\label{tab:signal-action}
\centering
\footnotesize
\begin{tabularx}{\columnwidth}{p{0.18\columnwidth} p{0.18\columnwidth} p{0.14\columnwidth} Y}
\toprule
\textbf{Signal} & \textbf{Action} & \textbf{Recovery} & \textbf{Interpretation} \\
\midrule
Idle compute proxy & Resize down & 60\% & Conservative efficiency recovery from a billing-side idle signal \\
Zombie resource & Remove resource & 100\% & Remove a resource when spend is effectively dormant \\
Tagging gap & Enforce tags & 0\% direct & Governance action that improves optimization readiness \\
Unattributed cost gap & Enforce tags & 0\% direct & Accountability repair rather than immediate cost reduction \\
Forecast risk & Review forecast risk & 0\% direct & Early warning rather than direct savings action \\
\bottomrule
\end{tabularx}
\end{table}

Current benchmark outputs show that the recommendation space is dominated by idle compute proxy actions. Query results report 48 medium-risk \texttt{resize\_down} recommendations accounting for \$1{,}545.74 in estimated savings, while governance-oriented \texttt{enforce\_tags} actions carry zero direct savings but remain important because unattributed NEC blocks accountable optimization.

\subsection{Forecasting and Lifecycle Validation}
The forecasting layer is intentionally lightweight. It estimates month-end NEC using a 70/30 blend of current run rate and trailing three-month average whenever sufficient history exists; otherwise it falls back to current run rate alone:
\begin{equation}
f = 0.7r + 0.3a
\end{equation}
or, when history is insufficient,
\begin{equation}
f = r.
\end{equation}
This model favors transparency over predictive complexity. The current benchmark backtest reports 68 evaluations with mean absolute error of \$7.86, mean absolute percentage error of 2.98\%, and a within-10-percent rate of 95.59\%.

The final layer treats recommendations as persisted operational objects rather than disposable dashboard rows. Lifecycle records store states such as \texttt{recommended}, \texttt{approved}, \texttt{implemented}, and \texttt{verified}, along with expected savings, realized savings, and derived accuracy. The current repository state should be described carefully: all 95 recommendations are still in the \texttt{recommended} state, with \$1{,}709.83 expected savings and \$0 realized savings. The lifecycle layer is therefore best presented as implemented validation instrumentation rather than mature realized-savings evidence.
